% Generated draft: requires mathematical and source review.
\documentclass[11pt]{article}
\usepackage[margin=25mm]{geometry}
\usepackage{fontspec}
\setmainfont{Latin Modern Roman}
\usepackage{amsmath,amssymb,mathtools,unicode-math}
\setmathfont{Latin Modern Math}
\usepackage{enumitem,needspace,xurl,hyperref}
\setlength{\parindent}{0pt}
\setlength{\parskip}{5pt}
\setlength{\emergencystretch}{3em}
\newcommand{\N}{\mathbb{N}}
\newcommand{\Z}{\mathbb{Z}}
\newcommand{\Q}{\mathbb{Q}}
\newcommand{\R}{\mathbb{R}}
\newcommand{\C}{\mathbb{C}}
\newcommand{\F}{\mathbb{F}}
\newcommand{\A}{\mathbb{A}}
\newcommand{\PP}{\mathbb{P}}
\newcommand{\E}{\mathbb{E}}
\newcommand{\eps}{\varepsilon}
\newcommand{\dd}{\,\mathrm{d}}
\begin{document}
\section*{UnsolvedMath ID: 30005120}
\textbf{Source problem code:} catalogue entry
\subsection{Definitions and mathematical statement}
Let \(n\to\infty\), and let \(G_{n,p}\) denote the random simple graph on \(\{1,\ldots,n\}\) in which each unordered pair of distinct vertices is an edge independently with probability \(p=p(n)\). Its chromatic number \(\chi(G_{n,p})\) is the least number of colours in a vertex colouring in which adjacent vertices receive different colours.

A deterministic concentration window consists of intervals \(I_n=[a_n,b_n]\), with endpoints allowed to depend on \(n\) and \(p\), such that
\[
\Pr\bigl(\chi(G_{n,p})\in I_n\bigr)\longrightarrow 1.
\]
Its length is \(b_n-a_n\). This gives a conventional precise meaning to concentration window; the supplied record does not specify the probability guarantee used in the referenced theorem.

The supplied question concerns the range
\[
n^{-1/2+\delta}\ll p(n)\leq 1-\delta,
\qquad
\frac{n^{-1/2+\delta}}{p(n)}\longrightarrow 0.
\]
The conventional interpretation takes \(\delta>0\) to be fixed independently of \(n\); a nonempty asymptotic range then requires \(0<\delta<1/2\). The record does not explicitly state the quantification of \(\delta\).

\textbf{Question.} For edge probabilities in this range, determine how close to optimal the supplied concentration-window upper bound
\[
\widetilde O\!\left(\frac{\sqrt{np}}{\log n}\right)
\]
is for \(\chi(G_{n,p})\). In particular, determine the order of the smallest achievable concentration windows and compare it with this bound. Here \(\widetilde O\) conventionally suppresses powers of logarithms, but the suppressed factors and their permitted parameter dependence are not specified in the record. The phrase ``close to optimal'' does not specify a comparison factor or whether matching lower bounds should hold for every admissible sequence \(p(n)\), only in particular regimes, or along subsequences. Consequently, the record supports this question but does not determine a unique quantified sharpness conjecture. The full statement of the referenced Theorem 17 is not supplied.

\textbf{Extraction qualification.} The original extracted text also appends a distinct ``Conjecture 16'': for every even \(k\geq 6\), if the solution space of
\[
a_1x_1+\cdots+a_kx_k=0,
\qquad
b_1x_1+\cdots+b_kx_k=0
\]
has dimension \(k-2\), then for sufficiently large \(q\) the pair of equations is ``uncommon''. The coefficient domain, the meaning and allowed values of \(q\), and the definition of ``uncommon'' are absent. This appended formulation is incompatible with the chromatic-number target and is retained here only as an unresolved extraction issue, not as an additional conjecture assigned to this permanent ID.
\subsection{Short English statement}
How small can an interval be while containing the chromatic number of \(G_{n,p}\) with probability tending to one, when \(n^{-1/2+\delta}\ll p\leq 1-\delta\)? The supplied question asks whether the stated \(\widetilde O(\sqrt{np}/\log n)\) upper bound is nearly optimal. Its logarithmic losses, probability guarantee, and intended lower-bound quantifiers require clarification. The supplied partially solved label is dated curation, not an independently verified conclusion.
\subsection{Sources}
\begin{itemize}
\item Supplied local UnsolvedMath registry record, permanent ID \(30005120\), title ``Sharp Concentration of Random-Graph Chromatic Number'', exact source problem code \texttt{catalogue entry}, in the Oberwolfach Reports Open Problems collection. No individual proposer is identified. The supplied proposed year is 2022.
\item Supplied report citation: \emph{Combinatorics, Probability and Computing} (2023), Oberwolfach Reports, DOI 10.4171/owr/2022/22, pp.\ 1165--1237. \url{https://doi.org/10.4171/owr/2022/22}. The original question refers to Theorem 17, whose statement is not included. The citation gives 2023, whereas the proposed year, year tag, and DOI contain 2022; this discrepancy is retained without resolution.
\item Supplied literature reference: Surya and Warnke, \emph{Electronic Journal of Combinatorics}. \url{https://doi.org/10.37236/11638}.
\item Supplied literature reference: Yan, ``very dense random graphs''. \url{https://arxiv.org/abs/2405.13914}.
\item Literature assessment checked at \texttt{2026-08-22}: the supplied status is ``partially solved''. The supplied assessment says that the stated upper concentration bounds were extended to sparse regimes, that later work establishes matching lower-order fluctuations in important dense regimes, and that a general best-possible characterization remains open. This is supplied dated curation derived from the OWR record, not independently verified here; its time-sensitive status and cited claims require independent verification.
\item The supplied statement status is \texttt{corrected\_verified}. Its verification note reports that the official Oberwolfach report was checked directly and that the clean statement restores definitions and removes extraction artifacts or source-dependent wording. This is a supplied verification claim, not a source visit or mathematical verification performed for this draft.
\end{itemize}
\end{document}
